% Antrag auf Ausstellung eines Bewohnerparkausweises (fiktives Buergeramt Musterstadt). German, 2 pages.
% Features: Eingangsstempel box, numbered sections with field boxes and checkboxes (page 1), signature block, field
% "Vom Amt auszufuellen"; page 2: Ausfuellhinweise, list of documents to attach, fees table, deadline, contact.
\documentclass[10pt,a4paper]{article}
\usepackage[T1]{fontenc}
\usepackage[utf8]{inputenc}
\usepackage[scaled=0.95]{helvet}
\renewcommand{\familydefault}{\sfdefault}
\usepackage[ngerman]{babel}
\usepackage[margin=1.8cm,top=1.8cm,bottom=1.8cm]{geometry}
\usepackage{xcolor}
\usepackage{tikz}
\usepackage{booktabs}
\usepackage{array}
\usepackage{tabularx}
\usepackage{enumitem}
\usepackage{fancyhdr}
\usepackage{lastpage}
\usepackage{microtype}

\definecolor{amt}{HTML}{8A1C1C}
\definecolor{feldgrau}{HTML}{6B6B6B}
\setlength{\parindent}{0pt}
\setlength{\parskip}{0.4em}
\setlength{\fboxsep}{3pt}

\pagestyle{fancy}
\fancyhf{}
\renewcommand{\headrulewidth}{0pt}
\fancyfoot[L]{\scriptsize Stadt Musterstadt \textperiodcentered{} Bürgeramt \textperiodcentered{} Formular BA-VP-07 \textperiodcentered{} Stand 10/2026 \textperiodcentered{} Fiktives Formular zu Testzwecken}
\fancyfoot[R]{\scriptsize Seite \thepage{} von \pageref{LastPage}}

% ---- Bausteine: Feld mit Beschriftung, Ankreuzfeld, Zeichenkasten, Abschnittsbalken ----------
\newcommand{\feld}[2]{\fbox{\begin{minipage}[t][0.95cm][t]{\dimexpr#1\linewidth-2\fboxsep-2\fboxrule\relax}%
  \scriptsize\textcolor{feldgrau}{#2}\end{minipage}}}
\newcommand{\kreuz}{\raisebox{-0.2ex}{\fbox{\rule{0pt}{1.35ex}\rule{1.35ex}{0pt}}}\,}
\newcommand{\zeichen}[1]{\foreach \i in {1,...,#1} {\fbox{\rule{0pt}{2.0ex}\rule{1.4ex}{0pt}}\hspace{-\fboxrule}}}
\newcommand{\sektion}[2]{\par\vspace{0.5em}\noindent\colorbox{amt}{\parbox{\dimexpr\linewidth-2\fboxsep}{\color{white}\bfseries
  #1\quad #2}}\par\vspace{0.3em}}

\begin{document}

% ---- Kopf ---------------------------------------------------------------------------------------
\noindent
\begin{minipage}[c]{0.09\linewidth}
\begin{tikzpicture}[scale=0.55]
  \draw[amt, fill=amt!15, line width=1.2pt] (0,0) -- (0,2) -- (2,2) -- (2,0) .. controls (2,-0.8) and (1,-1.2) .. (1,-1.4)
        .. controls (1,-1.2) and (0,-0.8) .. (0,0) -- cycle;
  \draw[amt, line width=1.2pt] (0.5,0.5) -- (1,1.5) -- (1.5,0.5) -- cycle;
\end{tikzpicture}
\end{minipage}\hfill
\begin{minipage}[c]{0.56\linewidth}
{\large\bfseries\color{amt} Stadt Musterstadt}\\
{\small Bürgeramt \textperiodcentered{} Abteilung Verkehr und Parken\\ Beispielplatz 1, 12345 Musterstadt}
\end{minipage}\hfill
\fbox{\begin{minipage}[t][2.0cm][t]{0.27\linewidth}\scriptsize\textcolor{feldgrau}{Eingangsstempel / Aktenzeichen\\
(vom Amt auszufüllen)}\end{minipage}}

\vspace{0.6em}
{\LARGE\bfseries Antrag auf Ausstellung eines Bewohnerparkausweises}\\[0.2em]
{\small Bitte in Druckbuchstaben ausfüllen und die Ausfüllhinweise auf Seite 2 beachten. Felder mit * sind Pflichtfelder.}

% ---- 1 ---------------------------------------------------------------------------------------------
\sektion{1}{Antragstellerin / Antragsteller}
\feld{0.49}{Name *}\hfill\feld{0.49}{Vorname(n) *}\\[0.2em]
\feld{0.30}{Geburtsdatum (TT.MM.JJJJ) *}\hfill\feld{0.33}{Telefon (tagsüber)}\hfill\feld{0.33}{E-Mail}\\[0.2em]
\feld{0.64}{Straße und Hausnummer *}\hfill\feld{0.16}{PLZ *}\hfill\feld{0.16}{Ort *}

% ---- 2 ---------------------------------------------------------------------------------------------
\sektion{2}{Fahrzeug}
Amtliches Kennzeichen *\quad \zeichen{3}\ --\ \zeichen{2}\ \zeichen{4}\quad (z.\,B. MU-AB 1234)\\[0.5em]
Fahrzeugart *\quad \kreuz Pkw \quad \kreuz Kraftrad \quad \kreuz Wohnmobil bis 3,5 t \quad \kreuz Anhänger\\[0.5em]
Die Halterin / der Halter des Fahrzeugs ist \quad \kreuz die Antragstellerin / der Antragsteller \quad \kreuz eine andere Person
oder Firma (Vollmacht und Halternachweis beifügen)\\[0.3em]
\feld{0.49}{Hersteller und Typ}\hfill\feld{0.49}{Name der Halterin / des Halters (falls abweichend)}

% ---- 3 ---------------------------------------------------------------------------------------------
\sektion{3}{Art des Antrags und Geltungsdauer}
\kreuz Erstausstellung \qquad \kreuz Verlängerung, bisherige Ausweisnummer: \fbox{\rule{0pt}{2.0ex}\rule{3.5cm}{0pt}} \\[0.3em]
\kreuz Ersatz bei Verlust oder Diebstahl \qquad \kreuz Änderung des Kennzeichens\\[0.5em]
Gewünschte Geltungsdauer *\quad \kreuz 1 Jahr \quad \kreuz 2 Jahre \qquad Gewünscht ab dem: \fbox{\rule{0pt}{2.0ex}\rule{2.8cm}{0pt}}

% ---- 4 ---------------------------------------------------------------------------------------------
\sektion{4}{Beigefügte Unterlagen}
\kreuz Kopie der Zulassungsbescheinigung Teil I (Fahrzeugschein) \qquad \kreuz Meldebescheinigung (nicht älter als 3 Monate)\\[0.3em]
\kreuz Vollmacht der Halterin / des Halters \qquad \kreuz Nachweis des Leasing- oder Firmenfahrzeugs \qquad \kreuz Anzeige bei Verlust

% ---- 5 ---------------------------------------------------------------------------------------------
\sektion{5}{Erklärung und Unterschrift}
Ich versichere, dass meine Angaben vollständig und richtig sind. Ich bin damit einverstanden, dass die Stadt
Musterstadt meine Daten zur Bearbeitung dieses Antrags speichert. Mir ist bekannt, dass der Ausweis nur für das
angegebene Fahrzeug im Bewohnerparkgebiet meiner Wohnanschrift gilt.\\[0.4em]
\feld{0.30}{Ort, Datum *}\hfill
\fbox{\begin{minipage}[t][1.3cm][t]{\dimexpr0.66\linewidth-2\fboxsep-2\fboxrule\relax}\scriptsize\textcolor{feldgrau}{Unterschrift der Antragstellerin / des Antragstellers *}\end{minipage}}

\vspace{0.6em}
\fcolorbox{feldgrau}{black!6}{\begin{minipage}{\dimexpr\linewidth-2\fboxsep-2\fboxrule\relax}
\small\textbf{Vom Amt auszufüllen}\\[0.2em]
Antrag eingegangen am: \rule{2.5cm}{0.4pt} \quad Gebühr bezahlt: \kreuz bar \ \kreuz Karte \ \kreuz Überweisung\\[0.5em]
Ausweisnummer: \rule{3cm}{0.4pt} \quad Gültig bis: \rule{2.5cm}{0.4pt} \quad Bearbeitet von: \rule{3cm}{0.4pt}
\end{minipage}}

\newpage
% ---- Seite 2 -----------------------------------------------------------------------------------------
\small
{\LARGE\bfseries\color{amt} Ausfüllhinweise}\\[0.1em]
{\small zum Antrag auf Ausstellung eines Bewohnerparkausweises}

\subsection*{Zu den einzelnen Feldern}
\begin{description}[leftmargin=2.4cm, labelwidth=2.2cm, style=nextline, itemsep=0.3em, font=\bfseries]
  \item[Zu 1] Tragen Sie Ihren Namen so ein, wie er im Personalausweis steht. Der Ausweis wird nur an Personen
        ausgestellt, die mit Hauptwohnsitz im Bewohnerparkgebiet gemeldet sind.
  \item[Zu 2] Geben Sie das Kennzeichen genau wie im Fahrzeugschein an. Pro Antrag kann nur ein Fahrzeug eingetragen
        werden. Gehört das Fahrzeug einer anderen Person oder einer Firma, benötigen Sie deren schriftliche Vollmacht.
  \item[Zu 3] Kreuzen Sie nur eine Art des Antrags an. Bei einer Verlängerung tragen Sie die Nummer des bisherigen Ausweises
        ein. Bei einem Verlust legen Sie bitte die Anzeige bei der Polizei bei.
  \item[Zu 4] Kreuzen Sie alle Unterlagen an, die Sie beifügen. Welche Unterlagen für Ihren Antrag nötig sind,
        steht in der Liste unten.
  \item[Zu 5] Der Antrag ist nur mit eigenhändiger Unterschrift gültig. Anträge per E-Mail ohne Unterschrift
        können nicht bearbeitet werden.
\end{description}

\subsection*{Diese Unterlagen müssen Sie beifügen}
\begin{enumerate}[leftmargin=1.6em, itemsep=0.2em]
  \item Kopie der Zulassungsbescheinigung Teil I (Fahrzeugschein) des Fahrzeugs.
  \item Meldebescheinigung, nicht älter als drei Monate (nur wenn Ihr Hauptwohnsitz nicht in Musterstadt gemeldet ist).
  \item Schriftliche Vollmacht und Kopie des Ausweises der Halterin / des Halters, wenn das Fahrzeug nicht auf Sie zugelassen ist.
  \item Bei Leasing- oder Firmenfahrzeugen: Bescheinigung der Leasinggesellschaft oder des Arbeitgebers über die dauerhafte Nutzung.
  \item Bei Ersatz wegen Verlust oder Diebstahl: Kopie der Anzeige bei der Polizei.
\end{enumerate}

\subsection*{Gebühren}
\begin{center}
\small
\begin{tabular}{@{}l r r@{}}
  \toprule
  \textbf{Leistung} & \textbf{Geltungsdauer} & \textbf{Gebühr} \\
  \midrule
  Erstausstellung oder Verlängerung & 1 Jahr & 15,00 Euro \\
  Erstausstellung oder Verlängerung & 2 Jahre & 30,00 Euro \\
  Ersatz bei Verlust oder Diebstahl & Restlaufzeit & 10,00 Euro \\
  Änderung des Kennzeichens & Restlaufzeit & 10,00 Euro \\
  \bottomrule
\end{tabular}
\end{center}
Die Gebühr wird bei Abholung des Ausweises im Bürgeramt fällig. Eine Rückerstattung ist nicht möglich.

\subsection*{Fristen und Bearbeitung}
\begin{itemize}[leftmargin=1.4em, itemsep=0.2em]
  \item Anträge auf \textbf{Verlängerung} müssen spätestens \textbf{4 Wochen vor Ablauf} des bisherigen Ausweises
        eingehen. Danach kann ein nahtloser Anschluss nicht garantiert werden.
  \item Die Bearbeitung dauert in der Regel \textbf{10 Werktage}, sobald alle Unterlagen vollständig vorliegen.
  \item Fehlen Unterlagen, fordert das Amt diese schriftlich nach. Werden sie nicht innerhalb von \textbf{14 Tagen}
        nachgereicht, wird der Antrag abgelehnt.
\end{itemize}

\subsection*{Abgabe und Kontakt}
Sie können den Antrag persönlich abgeben, per Post senden oder in den Briefkasten des Bürgeramts werfen.\\
\textbf{Bürgeramt Musterstadt, Abteilung Verkehr und Parken}, Beispielplatz 1, 12345 Musterstadt\\
Öffnungszeiten: Montag bis Freitag 8 bis 12 Uhr, Donnerstag zusätzlich 14 bis 18 Uhr \quad Telefon: 0000 555-0100

\subsection*{Datenschutz}
\small Ihre Daten werden nur zur Bearbeitung dieses Antrags verwendet und nach Ablauf der gesetzlichen Aufbewahrungsfrist
gelöscht. Weitere Informationen finden Sie in der Datenschutzerklärung der Stadt Musterstadt im Bürgeramt.

\end{document}
